\subsection*{1.3. 函子 $\Psi$}

\subsubsection*{1.3.1.}
设 $S$ 为 0.2.5 中的 Hensel 离散赋值环谱, 并设 $p:X\longrightarrow S$. 我们要定义一个从 $X$ 上的层到 1.2 所构造的拓扑斯 $X_s\times_sS$ 上的层的左正合函子 $\Psi$. 一般说来, 此函子不是某个拓扑斯态射的正像函子.

我们还将定义
\begin{equation*}
\Psi_\eta:(\text{层, 底为 }X_\eta)
\longrightarrow(\text{层, 底为 }X_s\times_s\eta).
\end{equation*}

\subsubsection*{1.3.2.}
置 $\bar X=X\times_S\bar S$. 有笛卡尔图
\begin{equation}\tag{1.3.2.1}
\begin{tikzcd}
X_{\bar s} \arrow[r, hook, "\bar i"] \arrow[d]
  & \bar X \arrow[d]
  & X_{\bar\eta} \arrow[l, hook', "\bar j"'] \arrow[d] \\
X_s \arrow[r, hook] & X & X_\eta \arrow[l, hook']
\end{tikzcd}
\end{equation}
位于图
\begin{center}
\begin{tikzcd}
\bar s \arrow[r, hook] \arrow[d]
  & \bar S \arrow[d]
  & \bar\eta \arrow[l, hook'] \arrow[d] \\
s \arrow[r, hook] & S & \eta \arrow[l, hook']
\end{tikzcd}
.\end{center}
之上. 设 $F$ 为 $X_\eta$ 上的层, 其在 $X_{\bar\eta}$ 上的逆像为 $F_{\bar\eta}$. 置
\begin{equation}\tag{1.3.2.2}
\Psi_\eta(F)=\bar i^*\bar j_*F_{\bar\eta}.
\end{equation}
通过结构转移, 这个层带有 $\operatorname{Gal}(\bar\eta/\eta)$ 的连续作用, 且与 $\operatorname{Gal}(\bar\eta/\eta)$ 在 $X_{\bar s}$ 上的作用相容. 函子
\begin{equation}\tag{1.3.2.3}
\Psi_\eta:(X_{\eta,\mathrm{et}})^{\sim}
\longrightarrow(X_s\times_s\eta)^{\sim}
\end{equation}
是左正合的.

\subsubsection*{1.3.3.}
设 $F$ 为 $X$ 上的层, $F_\eta$ 和 $F_s$ 为它在 $X_\eta$ 和 $X_s$ 上的限制. 置
\begin{equation}\tag{1.3.3.1}
(\Psi(F))_\eta=\Psi_\eta(F_\eta)
\end{equation}
及
\begin{equation}\tag{1.3.3.2}
(\Psi(F))_s=F_s.
\end{equation}
设 $\bar F$ 为 $F$ 在 $\bar X$ 上的逆像, $\varphi'$ 为伴随态射
\begin{equation*}
\varphi':\bar F\longrightarrow\bar j_*\bar j^*\bar F.
\end{equation*}
此态射诱导
\begin{equation*}
\varphi:F_{\bar s}=\bar i^*\bar F
\longrightarrow
\bar i^*\bar j_*\bar j^*\bar F
=(\Psi_\eta(F_\eta))_{\bar\eta}.
\end{equation*}
三元组
\begin{equation*}
\Psi(F)=((\Psi(F))_s,(\Psi(F))_\eta,\varphi)
\end{equation*}
是 $X_s\times_sS$ 上的层. 函子
\begin{equation}\tag{1.3.3.3}
\Psi:(X_{\mathrm{et}})^{\sim}
\longrightarrow(X_s\times_sS)^{\sim}
\end{equation}
是左正合的.

\subsubsection*{1.3.4.}
在不会引起混淆时, 我们记
\begin{equation*}
\Psi_\eta(F)=(\Psi(F))_\eta,
\qquad
\Psi_{\bar\eta}(F)=(\Psi(F))_{\bar\eta},
\qquad\text{及}\qquad
\Psi_{\bar\eta}(F)=(\Psi_\eta(F))_{\bar\eta}.
\end{equation*}

\subsubsection*{1.3.5.}
本节余下部分列举 $\Psi$ 的函子性质（参见 1.2.7--1.2.9）. 这些性质主要以其导出形式（2.1）发挥作用.

设 $f:X\longrightarrow X'$ 为 $S$-概形态射, 给出图
\begin{equation}\tag{1.3.5.1}
\begin{tikzcd}
X_{\bar s} \arrow[r, hook, "\bar i"] \arrow[d, "f"]
  & \bar X \arrow[d, "f"]
  & X_{\bar\eta} \arrow[l, hook', "\bar j"'] \arrow[d, "f"] \\
X'_{\bar s} \arrow[r, hook, "\bar i"]
  & \bar X' & X'_{\bar\eta} \arrow[l, hook', "\bar j"']
\end{tikzcd}
\end{equation}
（原印图误将上方中央的 $\bar X$ 漏印了上划线）. 箭头记号中的歧义简化了公式而不致引起混淆: 只有一种解释能使下面的复合有意义. 同样, 对 $X$ 和 $X'$ 的函子 (1.3.3.3), 我们都记为 $\Psi$.

下面对 $\Psi$ 描述的函子态射都有关于 $\Psi_\eta$ 的变体, 通常不再将其明确写出.

\subsubsection*{1.3.6.}
基变换态射 $\bar i^*f_*\longrightarrow f_*\bar i^*$ 诱导函子态射
\begin{equation}\tag{1.3.6.1}
\Psi f_*\longrightarrow f_*\Psi,
\end{equation}
当 $f$ 固有时它是同构（SGA~4, XII, 5.1(i)）.

关键情形是 $X'=S$. 此时得到态射
\begin{equation}\tag{1.3.6.2}
f_*\longrightarrow f_*\Psi
\end{equation}
其关键分量为
\begin{equation}\tag{1.3.6.3}
\Gamma(X_{\bar\eta},F)
\longrightarrow
\Gamma(X_{\bar s},\Psi_{\bar\eta}(F));
\end{equation}
当 $f$ 固有时这个态射是同构.

\subsubsection*{1.3.7.}
基变换态射 $f^*\bar j_*\longrightarrow\bar j_*f^*$ 诱导函子态射
\begin{equation}\tag{1.3.7.1}
f^*\Psi\longrightarrow\Psi f^*.
\end{equation}

\subsubsection*{1.3.8.}
假设 $f$ 拟有限. 函子 $f_!$ 与基变换交换; 态射 $f_!\bar j_*\longrightarrow\bar j_*f_!$ 诱导函子态射
\begin{equation}\tag{1.3.8.1}
f_!\Psi\longrightarrow\Psi f_!.
\end{equation}
当 $f$ 有限时, 它是 (1.3.6.1) 的逆.

\subsubsection*{1.3.9.}
对 $f^!$ 有对偶的情形（对阿贝尔层, 假设 $f$ 拟有限; 对带基点集合层, 假设该态射为浸入）. 这个函子与正像交换, 而态射 $\bar i^*f^!\longrightarrow f^!\bar i^*$ 诱导函子态射
\begin{equation}\tag{1.3.9.1}
\Psi f^!\longrightarrow f^!\Psi.
\end{equation}
当 $f$ 平展时, 它是 (1.3.7.1) 的逆.

\subsubsection*{1.3.10.}
最后, 一个基变换 $S'\xrightarrow{f}S$ 给出基变换态射
\begin{equation}\tag{1.3.10.1}
f^*\Psi\longrightarrow\Psi f^*.
\end{equation}
